\section{Tujuan Buku Ajar}
\label{sec:tujuan-buku}

Buku ajar ini disusun untuk satu tujuan yang jelas: membuat mahasiswa
mampu \emph{bekerja} pada level instruksi mesin, bukan sekadar
mengetahui bahwa level itu ada. Mata kuliah Pemrograman Bahasa
Tingkat Rendah bukanlah mata kuliah hafalan mnemonic. Yang dinilai
bukan seberapa banyak instruksi yang dapat diingat, melainkan
seberapa mampu mahasiswa membaca sebuah persoalan, memilih
instruksi yang tepat, menuliskannya menjadi program yang benar, dan
memperbaiki program tersebut ketika hasilnya salah.

\subsection{Learning Outcomes Buku}

Pada akhir seluruh bab, mahasiswa diharapkan telah mencapai empat
kemampuan utama berikut.

\begin{enumerate}
    \item \textbf{Menjelaskan cara kerja komputer pada level mesin.}
          Mencakup arsitektur internal 8086, register, segmentasi
          memori, dan siklus instruksi.

    \item \textbf{Membaca dan menjelaskan hubungan instruksi mesin
          dengan perangkat keras.} Mencakup format instruksi, mode
          pengalamatan, direktif assembler, dan mekanisme interupsi.

    \item \textbf{Menulis program assembly.} Mencakup operasi
          transfer data, aritmatika, logika, struktur kontrol,
          prosedur, dan operasi masukan/keluaran.

    \item \textbf{Melakukan penelusuran dan perbaikan program pada
          level instruksi.} Mencakup penggunaan \eng{debugger},
          analisis memori dan stack, pembacaan hasil
          \eng{disassembly}, serta optimasi.
\end{enumerate}

Keempat kemampuan tersebut bukan daftar topik yang berdiri sendiri.
Kemampuan kedua bergantung pada kemampuan pertama; kemampuan ketiga
memerlukan kedua kemampuan sebelumnya; dan kemampuan keempat baru
bermakna setelah mahasiswa mampu menulis program. Karena itu urutan
bab dalam buku ini disusun mengikuti ketergantungan tersebut dan
tidak disarankan dilompati.

\subsection{Kompetensi Akhir yang Diukur}

Tabel~\ref{tab:outcome-buku} memetakan keempat kemampuan di atas
kepada capaian pembelajaran yang ditetapkan pada mata kuliah.

\begin{table}[H]
    \centering
    \caption{Pemetaan kemampuan akhir buku terhadap CPMK dan CPL.}
    \label{tab:outcome-buku}
    \small
    \begin{tabular}{@{}p{4.6cm}p{1.9cm}p{1.9cm}p{4.3cm}@{}}
        \toprule
        \textbf{Kemampuan akhir buku} & \textbf{CPMK} &
        \textbf{CPL} & \textbf{Bab utama} \\
        \midrule
        Menjelaskan arsitektur register dan manajemen memori &
        \CPMK{01} & \CPL{02} & Bab~IV--VI \\
        \addlinespace
        Memahami hubungan instruksi mesin dengan perangkat keras &
        \CPMK{02} & \CPL{02} & Bab~VII--IX, XIV \\
        \addlinespace
        Menulis kode assembly untuk operasi aritmatika dan logika &
        \CPMK{03} & \CPL{03} & Bab~X--XIII \\
        \addlinespace
        Melakukan \eng{debugging} pada level instruksi memori &
        \CPMK{04} & \CPL{03} & Bab~XV--XVI \\
        \bottomrule
    \end{tabular}
\end{table}

\subsection{Apa yang Tidak Dicakup Buku Ini}

Agar batas pembahasan jelas, buku ini \emph{tidak} membahas:
pemrograman 32-bit maupun 64-bit secara rinci, sintaks NASM, arsitektur
prosesor selain keluarga 8086, dan pemrograman sistem operasi modern.
Pembahasan sengaja dibatasi pada 8086 karena prosesor inilah yang
paling bersih untuk mempelajari konsep dasar: jumlah register
terbatas, tanpa mode terproteksi, dan seluruh perilakunya dapat
diamati dengan emulator sederhana. Konsep yang dipelajari di sini
berlaku langsung pada keluarga x86 modern.

\begin{infobox}
Bila Anda menemukan istilah yang belum dikenal, Lampiran~C memuat
glosarium Indonesia--Inggris. Setiap kemunculan pertama istilah asing
pada bab-bab berikutnya juga ditandai dengan huruf miring, misalnya
\eng{instruction queue}.
\end{infobox}
